
\documentclass[../../main.tex]{subfiles}
% \documentclass[a4paper,10pt]{ltjsarticle}
\begin{document}
% \documentclass[a4paper,12pt]{article}
% \usepackage[utf8]{inputenc}
% \usepackage[japanese]{babel}
% \usepackage{xeCJK}
% \usepackage{geometry}
% \usepackage{array}
% \usepackage{booktabs}
% \usepackage{fancybox}

% % ページ設定
% \geometry{top=25mm, bottom=25mm, left=20mm, right=20mm}
% \pagestyle{empty}

% % フォント設定
% \setCJKmainfont{Noto Sans CJK JP}

% \begin{document}

% 年度の大きな表示
\begin{center}
  {\fontsize{48pt}{58pt}\selectfont \textbf{2001}\normalsize\textbf{年度}}
\end{center}

\vspace{10mm}

% 本文


\vspace{15mm}

% 出題テーマと難易度の表
\noindent\textbf{出題テーマと難易度}

\vspace{5mm}

\begin{center}
  \shadowbox{%
    \begin{tabular}{|c|p{8cm}|c|}
      \hline
      \textbf{問題番号} & \textbf{テーマ} & \textbf{難易度} \\
      \hline
      第1問           & 三角関数，数列      & 普            \\
      \hline
      第2問           & 平面図形         & 難            \\
      \hline
    \end{tabular}%
  }
\end{center}

\vspace{15mm}

% 解答
\noindent\textbf{解答}

\vspace{5mm}

\begin{center}
  \shadowbox{%
    \renewcommand{\arraystretch}{1.6}
    \begin{tabular}{|c|c|p{8cm}|}
      \hline
      \textbf{問題番号}        & \textbf{小問} & \textbf{答え}                                           \\
      \hline
      \multirow{4}{*}{第1問} & (1)         & 証明問題                                                  \\
      \cline{2-3}
                           & (2)         & 証明問題                                                  \\
      \cline{2-3}
                           & (3)         & 証明問題                                                  \\
      \cline{2-3}
                           & (4)         & 証明問題                                                  \\
      \hline
      \multirow{2}{*}{第2問} & (1)         & $\displaystyle    \begin{dcases}
                                                                 a = \frac{\sqrt{2}(\sqrt{3}+1)}{2} \\
                                                                 b = \sqrt{3}+1
                                                               \end{dcases}$ \\
      \cline{2-3}
                           & (2)         & $\displaystyle R = \frac{\sqrt{2}(\sqrt{3}+1)}{2}$    \\
      \hline
    \end{tabular}%
    \renewcommand{\arraystretch}{1.0}
  }
\end{center}

% 解答時間
\noindent\textbf{試験形態}

\vspace{5mm}

\begin{center}
  \shadowbox{%
    \begin{tabular}{p{8cm}}
      2問90分
    \end{tabular}%
  }
\end{center}


\end{document}